\section*{الفصل \ref{ch:b3:bioinorganic} --- الكيمياء غير العضوية الحيوية}

\begin{solution}{exo:b3:bioinorganic:1}
ينتقل $\log(\theta/(1 - \theta))$ من $\log 0.25 = -0.602$ إلى $\log 4 = +0.602$ بينما يرتفع $\log p$ بمقدار
$\log(5.9/2.1) = 0.449$: $n = 1.204/0.449 = 2.7$.
\end{solution}

\begin{solution}{exo:b3:bioinorganic:2}
$\ce{Mn^2+} < \ce{Fe^2+} < \ce{Co^2+} < \ce{Ni^2+} < \ce{Cu^2+} > \ce{Zn^2+}$. ليس للزنك ($d^{10}$) استقرار مجال ربائط ولا
تشوه يان--تيلر: فيعود إلى ما دون النحاس.
\end{solution}

\begin{solution}{exo:b3:bioinorganic:3}
$\ce{Ca^2+}$ و$\ce{Fe^3+}$، قاسيان: الكربوكسيلات (والفينولات للحديد). $\ce{Cu+}$ و$\ce{Hg^2+}$، لينان:
ثيولات السيستئين. $\ce{Zn^2+}$، حدّي: الهستيدين والسيستئين.
\end{solution}

\begin{solution}{exo:b3:bioinorganic:4}
$\ce{2H2O -> O2 + 4H+ + 4e-}$: أربعة إلكترونات، وأربعة فوتونات لكل $\ce{O2}$؛ وأربعة مولات من الفوتونات لكل
مول من $\ce{O2}$.
\end{solution}

\begin{solution}{exo:b3:bioinorganic:5}
الميوغلوبين: $\theta = 13/13.37 = 0.972$ و$5/5.37 = 0.931$: يتخلى عن $(0.972 - 0.931)/0.972 = 4$\,\% من
أكسجينه. الهيموغلوبين: $\theta = 0.972$ و$0.724$: يتخلى عن 26\,\%.
\end{solution}

\begin{solution}{exo:b3:bioinorganic:6}
$[\mathrm{HbCO}]/[\mathrm{HbO_2}] = 230 \times \num{50e-6}/0.2095 = 0.055$؛ والنسبة $0.055/1.055 = 5.2$\,\%.
\end{solution}

\begin{solution}{exo:b3:bioinorganic:7}
الكبريتيدات $-2$ لكل منها، والسيستئينات $-1$ لكل منها. $\ce{[Fe2S2(SR)4]^2-}$: مجموع الحديد $-2 + 4 + 4 = +6$: ذرتا
Fe(III). $\ce{[Fe2S2(SR)4]^3-}$: $+5$، ذرة Fe(III) وذرة Fe(II) (تكافؤ مختلط). $\ce{[Fe4S4(SR)4]^2-}$:
$-2 + 8 + 4 = +10$، ذرتا Fe(III) وذرتا Fe(II)، أي $+2.5$ لكل حديد في المتوسط.
\end{solution}

\begin{solution}{exo:b3:bioinorganic:8}
$\ce{[PtCl2(NH3)2] + H2O -> [PtCl(H2O)(NH3)2]+ + Cl-}$. فتركيز الكلوريد المرتفع في الدم
يدفع التوازن عائدًا؛ وداخل الخلايا يكون الكلوريد أدنى بكثير، فيتكوّن المعقد
المائي. ويرتبط البلاتين بالذرة N7 في الغوانين. وفي المتماكب \textit{trans} يكون الموقعان
السريعا الاستبدال متقابلين فلا يستطيعان بلوغ قاعدتين متجاورتين على
سلسلة واحدة.
\end{solution}

\begin{solution}{exo:b3:bioinorganic:9}
$1/T_1 = 1/\qty{1.2}{s} + \qty{4.0}{L.mmol^{-1}.s^{-1}} \times \qty{0.10}{mmol/L} = 0.833 + 0.400 = \qty{1.23}{s^{-1}}$: $T_1 = \qty{0.81}{s}$.
\end{solution}

\begin{solution}{exo:b3:bioinorganic:10}
$K = K_{\mathrm{PbL}}/K_{\mathrm{CaL}} = 10^{18.0 - 10.7} = 10^{7.3} = \num{2e7}$: فالرصاص يزيح الكالسيوم كليًا. والربيطة
الحرة كانت سترتبط أيضًا بكالسيوم الدم وتخفض تركيزه
تخفيضًا خطرًا؛ أما معقد الكالسيوم فيتبادل مع الرصاص ويترك مستوى الكالسيوم
دون تغيير.
\end{solution}

\begin{solution}{exo:b3:bioinorganic:11}
بعيدًا تحت $K_M$ يكون التفاعل من الرتبة الأولى: $k' = (k_{\mathrm{cat}}/K_M)[\mathrm E]$، $t_{1/2} = \ln 2/k'$. أنهيدراز
الكربونيك: $k' = \qty{100}{s^{-1}}$، $t_{1/2} = \qty{6.9}{ms}$. \omterm{def:b3:complex-kinetics:enzyme}{والإنزيم} النموذجي: $k' = \qty{0.10}{s^{-1}}$، $t_{1/2} = \qty{6.9}{s}$، أي أطول بألف
مرة، وهو أبطأ من أن يناسب الثانية التي تقضيها كرية حمراء في شعيرة رئوية.
\end{solution}

\begin{solution}{exo:b3:bioinorganic:12}
نصف المواقع مع CO يعني $[\mathrm{HbCO}] = [\mathrm{HbO_2}]$: $p(\ce{CO}) = p(\ce{O2})/M = 0.2095/230 = \qty{9.1e-4}{bar}$، أي نحو
\qty{910}{ppm}.
\end{solution}

\begin{solution}{pb:b3:bioinorganic:1}
\textbf{1.} $\theta = 13^{2.7}/(3.5^{2.7} + 13^{2.7}) = 0.972$.

\textbf{2.} $\theta(\qty{5.0}{kPa}) = 0.724$.

\textbf{3.} $5.0/(0.37 + 5.0) = 0.931$.

\textbf{4.} قيمة $p_{50}$ له أدنى بعشر مرات: فعند ضغوط العضلة العاملة يكون أكثر
إشباعًا بكثير من الهيموغلوبين عند الضغط نفسه، فينتقل الأكسجين من
الهيموغلوبين إلى الميوغلوبين.

\textbf{5.} $n = 1$.

\textbf{6.} الارتباط عند \omterm{def:b3:bioinorganic:haem}{هيم} واحد ينقل حديده إلى المستوي ويسحب الهستيدين 
وحلزونه؛ فتنزاح الوحدات معًا إلى صورة أعلى ألفة إزاء
المواقع الأخرى.

\textbf{7.} $\qty{150}{g/L}/\qty{64500}{g/mol} = \qty{2.33e-3}{mol/L}$.

\textbf{8.} $4 \times \num{2.33e-3} = \qty{9.30e-3}{mol/L}$.

\textbf{9.} $\qty{9.30}{mmol/L} \times 0.972 = \qty{9.04}{mmol}$.

\textbf{10.} $\qty{9.30}{mmol/L} \times (0.972 - 0.724) = \qty{2.31}{mmol}$.

\textbf{11.} $\qty{2.31e-3}{mol} \times \qty{22.41}{L/mol} = \qty{52}{mL}$.

\textbf{12.} $0.248/0.972 = 26$\,\%.

\textbf{13.} $5.0^{2.7}/(4.0^{2.7} + 5.0^{2.7}) = 0.646$.

\textbf{14.} $\qty{9.30}{mmol/L} \times (0.972 - 0.646) = \qty{3.03}{mmol}$.

\textbf{15.} $3.03/2.31 = 1.31$: أي أكثر بثلاثين في المئة.

\textbf{16.} ثنائي أكسيد الكربون الناتج عن التنفس، الذي يميّهه أنهيدراز الكربونيك إلى حمض
الكربونيك، وحمض اللبن في الجهد الشديد.

\textbf{17.} $\qty{3.03}{mmol/L} \times \qty{5.0}{L/min} = \qty{15.1}{mmol/min}$.

\textbf{18.} $\qty{15.1e-3}{mol/min} \times \qty{22.41}{L/mol} = \qty{0.34}{L/min}$.

\textbf{19.} $230 \times \num{100e-6}/0.2095 = 0.110$.

\textbf{20.} $0.110/1.110 = 0.099$: أي نحو 10\,\% من المواقع.

\textbf{21.} المواقع التي لا تزال تحمل الأكسجين، في جزيء يحمل CO أيضًا، تكون في
الصورة العالية الألفة: فينزاح المنحنى نحو اليسار ولا تحرر إلا قليلًا من
أكسجينها في الأنسجة.

\textbf{22.} بحسب علاقة التنافس تنخفض النسبة $[\mathrm{HbCO}]/[\mathrm{HbO_2}]$ بالتناسب مع $p(\ce{O2})$:
فاستنشاق الأكسجين النقي، بضغط يقارب خمسة أضعاف ضغطه في الهواء، يزيح أحادي
أكسيد الكربون أسرع.

\textbf{23.} بإعاقة الارتباط الخطي للجزيء CO وتكوين رابطة هيدروجينية مع $\ce{O2}$، يُبقي $M$ في
حدود المئات بدل القيمة الأكبر بكثير \omterm{def:b3:bioinorganic:haem}{للهيم} الحر؛ وإلا لحجب حتى
أحادي أكسيد الكربون الذي يصنعه الجسم جزءًا كبيرًا من الهيموغلوبين.

\textbf{24.} في الراحة يوصل الدم نحو \qty{0.34}{L} من الأكسجين في الدقيقة.
\end{solution}
